\documentclass[../report.tex]{subfiles}
\begin{document}

\section{Datapath}

Trong CPU đơn chu kỳ, toàn bộ đường từ thanh ghi \sig{PC} qua các module đến register file và RAM là mạch tổ hợp.
Mỗi lệnh xong trong một chu kỳ: ở cạnh lên của clock, \sig{PC}, thanh ghi đích và RAM cùng cập nhật. Hình~\ref{fig:sc}
là mạch top-level \sig{RV32I\_SingleCycle}.

\begin{figure}[H]
\centering
\includegraphics[width=\textwidth]{RV32I_SingleCycle.png}
\caption{Mạch top-level của CPU đơn chu kỳ trong Logisim}\label{fig:sc}
\end{figure}

Một lệnh đi qua năm bước, tất cả trong cùng một chu kỳ:

\begin{enumerate}
  \item Nạp lệnh: \sig{PC} đưa vào \sig{IMEM} (\sig{PC[11:2]}) cho ra \sig{instr}; song song đó bộ cộng tính PC + 4.
  \item Giải mã: \sig{Control} sinh tín hiệu điều khiển, \sig{RegisterFile} đọc rs1 và rs2, \sig{ImmGen} tạo
        immediate.
  \item Thực thi: A = rs1, PC hoặc 0 theo \sig{ASel}; B = rs2 hoặc immediate theo \sig{BSel}; \sig{ALU} tính,
        \sig{BranchUnit} so sánh rs1 với rs2.
  \item Truy cập bộ nhớ: địa chỉ là kết quả ALU. \sig{StoreUnit} ghép dữ liệu ghi, RAM ghi ở cạnh lên nếu
        \sig{RamWE} = 1, \sig{LoadUnit} tách dữ liệu đọc.
  \item Ghi lại: \sig{WBSel} chọn kết quả ALU, dữ liệu load hoặc PC + 4; register file ghi vào rd ở cạnh lên.
\end{enumerate}

PC kế tiếp được chọn bởi \sig{PCSel} = \sig{Taken} $\lor$ \sig{Jal} $\lor$ \sig{Jalr}: bằng 1 thì PC nhận kết quả
ALU AND \texttt{0xFFFFFFFE}, bằng 0 thì nhận PC + 4. Đặc tả JALR yêu cầu xóa bit 0 của đích nhảy \cite{riscv-spec};
với branch và JAL bit này vốn bằng 0, nên một mạch dùng chung cho cả ba. Thanh ghi \sig{PC} có enable = NOT
\sig{Halt}, nên khi gặp EBREAK, CPU giữ nguyên trạng thái và chân \sig{halt} lên 1.

\section{Ví dụ: lệnh \texttt{lw t4, 8(sp)}}

Lệnh này có mã máy \texttt{0x00812E83}: immediate = 8, rs1 = x2, funct3 = 010, rd = x29, opcode = 0000011 (LOAD).
\sig{Control} bật nhóm LOAD và cho \sig{RegWrite} = 1, \sig{ASel} = 0, \sig{BSel} = 1, \sig{ImmSel} = 0,
\sig{MemRead} = 1, \sig{MemWrite} = 0, \sig{WBSel} = 1. ALU cộng sp với 8 ra địa chỉ, \sig{AddrDecode} thấy
\sig{IsIO} = 0 nên đọc RAM, \sig{LoadUnit} chọn cả word theo funct3 = 010, MUX ghi lại chọn dữ liệu bộ nhớ, và ở
cạnh lên t4 nhận giá trị. \sig{PCSel} = 0 nên PC nhận PC + 4.

\section{Đường găng và giới hạn của thiết kế đơn chu kỳ}\label{sec:critical}

Chu kỳ clock phải đủ dài cho đường tổ hợp dài nhất. Lệnh load đi qua nhiều khối nhất: \sig{IMEM}, \sig{Control}
và \sig{RegisterFile}, \sig{ALU}, \sig{DMEM}, \sig{LoadUnit}, MUX ghi lại, rồi đến đầu vào register file. Các lệnh
ngắn hơn như \texttt{add} không dùng RAM nhưng vẫn phải chờ hết chu kỳ đó.

Logisim không mô phỏng trễ cổng, nên báo cáo dùng giả định trễ từng bước theo Patterson và Hennessy
\cite[ch.~4]{patterson}: nạp lệnh 200~ps, đọc thanh ghi 100~ps, ALU 200~ps, truy cập bộ nhớ dữ liệu 200~ps, ghi thanh
ghi 100~ps. Với giả định này, chu kỳ đơn chu kỳ là \resTSingle~ps. Đây là con số giả định để so sánh, không phải
số đo trên mạch. Chương~\ref{ch:pipe} chia đường này thành năm tầng để rút ngắn chu kỳ clock.

\end{document}
